\section{From the literal table to the consistency statement}
\label{sec:endpoint}

We now give the precise composition proving Theorem~\ref{thm:main}.
It is useful to separate a theorem of ordinary mathematics from the
additional work needed to serialize that theorem in a proof assistant
or in an arithmetized theory.

\subsection{The calibrated premise is a concrete source theorem}

The pinned ProveIt declaration
\file{CalibratedRotation240.halts_iff_zfc_refutation} states
\[
 \RHalt(P_0,(\epsilon,R_*))
 \leftrightarrow
 \operatorname{BProv}(\mathrm{ZFCax\_s},[],\bot).
\tag{C}
\]
Here $\operatorname{BProv}$ denotes existence of a finite first-order
derivation in the explicitly supplied ZFC calculus. The new source-data
check identifies $P_0$ and $R_*$ with its actual definitions. The
earlier ordinary calibration includes the proof-program enumeration,
its compact logical basis and native CP/GIC theory, the quotient
interpretation of ZF, and the comparison of ZF and ZFC refutations.
It is not merely an appeal to the informal word ``equiconsistent.\cite{calibration,upstream-settheory}''

The source and binary theorem declarations read at the current checkpoint
are retained under \file{reference/current/}. The current
\file{BinaryRotation240.nonhalting_iff_zfc_consistency} is an external
Lean endpoint for the original 240 table. We rely on the mathematical
content of this development and its cited calibration, not on a claimed
general soundness theorem for all of Lean's type theory inside ZFC.
The underlying arguments about finite syntax, models of first-order
theories, quotients, and definable submodels are formalizable in ZFC.

The existence of a finite $\operatorname{BProv}$ derivation is equivalent
to $\exists p\,\Proof_{\mathrm{ZFC}}(p,\bot)$ by effective encoding
and decoding of finite proof trees. Other customary effective proof
codings give the same consistency statement through finite syntactic
translations. No truth predicate for the universe of all sets is
required by this change of proof coding.

\subsection{Proof of the main theorem}

By the complete blank prefix, $M_{238}$ reaches a valid frame for
$(P_1,\epsilon,R_*)$. By the unbounded runtime proof, every source step
corresponds to a finite positive binary segment and has the exact
claimed successor configuration. A source halt gives a finite binary
halt. Conversely an infinite nonhalting source execution supplies
infinitely many positive binary segments; it cannot be represented
by a finite binary prefix or be trapped in one instruction. A premature
binary halt is excluded by the complete segment computations. Therefore
\[
       \Halt(M_{238},0^{\mathbb Z})
       \leftrightarrow \RHalt(P_1,(\epsilon,R_*)).
\tag{R}
\]

The exact source control certificate gives
\[
       \RHalt(P_1,(\epsilon,R_*))
       \leftrightarrow \RHalt(P_0,(\epsilon,R_*)).
\tag{S}
\]
Composing (R), (S), and (C), followed by finite-proof coding, gives
\[
       \Halt(M_{238},0^{\mathbb Z})
       \leftrightarrow \exists p\,\Proof_{\mathrm{ZFC}}(p,\bot).
\]
Negating both sides yields the stated consistency equivalence. The
canonical state renumbering preserves each binary transition, so the
same theorem holds for the canonical table in the appendix.

All new side conditions are finite facts about exact tables and source
data, proved by explicit calculations. All new unbounded arguments are
inductions on natural numbers, finite lists, or source execution lengths.
Together with the inherited calibration mathematics, they give the
ordinary ZFC proof claimed here. \hfill$\square$

The separately audited local route supplies a useful cross-check:
boundary239 already has the calibrated source/runtime interface, and
the new selector guard plus the delayed-erasure blocks prove that
final 238 halts exactly when boundary239 does. The direct runtime route
above does not assume that selector guard, so a mistake in its extraction
would not silently become a premise of the direct final-table proof.

\subsection{The fresh guard's exact scope}

The guard replay examines the complete boundary239 startup and all 8,386
unbounded cases for $P_1$, using 2,247,513 symbolic rules. There are 96
symbolic visits to inherited state 169, of which 32 read a one. At every
such positive visit the exposed one-run is
\[
                         G+K,\qquad23\le K\le59.
\]
Only the gap coefficient is nonzero; all other variable coefficients
vanish in these run lengths. Thus the required neighboring one is
forced for every $G\ge0$. All 32 positive visits occur at increment
sites for logical register 4. The literal predecessor startup has 18
visits to 169, including six that read a one, and each is checked directly.
The four list schemas do not visit 169, and 169 has no self-loop, so no
accelerated rule can hide an uncovered visit.

\subsection{Formalization and independence distinctions}

The result is an ordinary proof formalizable in ZFC, with reproducible
executable certificates. It is not a newly compiled external Lean/Rocq
theorem for this exact 238 table. It also does not claim that the complete
equivalence has been proved inside PA. The existing repository contains
separate ongoing work on arithmetized endpoints; that work is not
silently completed by a Python certificate or by an external theorem.

For a future exact-table Lean endpoint, the local route can port the
finite block relations and the fresh guard into the existing machine
semantics. The direct route can instantiate the existing runtime
framework for this table and certify its exact source coverage and
409,835-step frame. Either must refer to the delivered table's literal
identity, not just a generic machine with the same state count.
Both routes must also check the $P_1\leftrightarrow P_0$ source bridge
in the kernel. The local route needs a checked boundary predecessor's
runtime and startup; the inherited closed binary endpoint belongs to
rotation240.

If ZFC is consistent, the machine does not halt. Under the usual
derivability conditions, a ZFC proof of this nonhalting statement would
give a ZFC proof of its own consistency, contrary to the second
incompleteness theorem. Bare consistency alone is not used to rule out
a ZFC proof of halting. A stronger hypothesis, such as arithmetical
$\Sigma_1$-soundness, supplies the familiar two-sided independence
conclusion. No minimum possible state count is established.
